\documentclass[aps,pra,twocolumn,superscriptaddress,floatfix,10pt]{revtex4-2}

\usepackage{amsmath,amssymb}
\usepackage{graphicx}
\usepackage{bm}

% ---- shared macros (match main text) ----
\newcommand{\CV}{\mathrm{CV}}
\newcommand{\Itr}{I_{\mathrm{tr}}}
\newcommand{\Ival}{I_{\mathrm{val}}}
\newcommand{\thetaz}{\theta_0}
\newcommand{\thetaK}{\theta_K}
\newcommand{\Ltr}{L_{\mathrm{tr}}}
\newcommand{\Lval}{L_{\mathrm{val}}}
\newcommand{\gtr}{g_{\mathrm{tr}}}
\newcommand{\gval}{g_{\mathrm{val}}}
\newcommand{\Vtr}{V_{\mathrm{tr}}}

\renewcommand{\thesection}{S\arabic{section}}
\renewcommand{\thetable}{S\Roman{table}}
\renewcommand{\thefigure}{S\arabic{figure}}
\renewcommand{\theequation}{S\arabic{equation}}

\begin{document}

\title{Supplemental Material:\\
Gradient concentration constrains trajectory diagnostics
in variational quantum meta-learning}

\author{Sang Ho Oh}
\affiliation{Department of Computer Engineering and Artificial Intelligence,
Pukyong National University, Busan 48513, Korea}

\maketitle

This Supplemental Material collects the supporting analyses referenced in the main text: the
concentration of the identity prefactor (Sec.~\ref{sm:prefactor}), the one-step versus multi-step
identity check (Sec.~\ref{sm:kstep}), the association of alignment with the held-out loss rather than
the gap (Sec.~\ref{sm:heldout}), the near-zero-mean-transfer character of the task family
(Sec.~\ref{sm:transfer}), and four robustness checks---a term-type-stratified split, additional
Hamiltonian families and entangler topologies, a finite-shot variance proxy, and a random
initialization (Sec.~\ref{sm:robust}); and a synthetic positive control showing that the joint
control passes a genuine alignment--gap signal (Sec.~\ref{sm:positive}). All quantities use the definitions, estimators, and
hierarchical-bootstrap procedure of the main text. Rank correlations are Spearman coefficients; joint
controls are partial rank correlations $\rho(A,G\,|\,\Ival,\Lval(\thetaz))$ unless stated otherwise.
The underlying per-task tables are provided in the accompanying data repository.

\section{Concentration of the identity prefactor}
\label{sm:prefactor}

The first-order identity of the main text reads $\Ival\approx\alpha\,A\,q$ with prefactor
$q=\|\gtr\|\|\gval\|$. The controlling role of gradient concentration is a statement about the
relative spread of $q$, not of $\|\gtr\|$ alone. Table~\ref{tab:cv} reports the pooled coefficient of
variation of $\|\gtr\|$, $\|\gval\|$, and $q$ across tasks at each system size. Both gradient norms
concentrate with scale, and the product $q$---being the product of two fluctuating norms---has a
larger but likewise decreasing coefficient of variation. This is the regime in which the ranking of
$\Ival$ collapses onto that of $A$.

\begin{table}[b]
\caption{\label{tab:cv}%
Pooled coefficient of variation of the two gradient norms and of the identity prefactor
$q=\|\gtr\|\|\gval\|$, as a function of system size $n$ (pooled over depth $L$ and seeds).}
\begin{ruledtabular}
\begin{tabular}{cccc}
$n$ & $\CV(\|\gtr\|)$ & $\CV(\|\gval\|)$ & $\CV(q)$ \\
\colrule
$4$  & $0.336$ & $0.361$ & $0.551$ \\
$6$  & $0.285$ & $0.282$ & $0.466$ \\
$8$  & $0.246$ & $0.261$ & $0.406$ \\
$10$ & $0.264$ & $0.261$ & $0.451$ \\
\end{tabular}
\end{ruledtabular}
\end{table}

\section{One-step versus multi-step identity}
\label{sm:kstep}

The identity is a first-order (tangent) expression; it is at its most accurate in the one-step limit.
Table~\ref{tab:kstep} compares the rank correlation between the alignment $A$ and the validation
improvement measured after a single inner step ($\Ival^{(1)}$) with that measured after the full
$K=10$ inner steps ($\Ival^{(K)}$). At $K=1$ the correlation is near-perfect and essentially flat in
$n$ ($0.95$--$0.97$), confirming that the tangent expression captures the one-step improvement
directly. At $K=10$ the correlation is lower and rises with scale ($0.60$ to $0.85$): as gradients
concentrate and adaptation steps shrink, the full trajectory moves closer to its initial tangent, and
the identity increasingly predicts the multi-step improvement as well. The tightening reported in the
main text is therefore the convergence of the $K$-step improvement toward the one-step tangent, not a
change in the identity itself.

\begin{table}[t]
\caption{\label{tab:kstep}%
Rank correlation of the support--validation alignment $A$ with the validation improvement after one
inner step ($\Ival^{(1)}$) and after $K=10$ inner steps ($\Ival^{(K)}$), pooled over depth and seeds.}
\begin{ruledtabular}
\begin{tabular}{ccc}
$n$ & $\rho(A,\Ival^{(1)})$ & $\rho(A,\Ival^{(K)})$ \\
\colrule
$4$  & $+0.95$ & $+0.60$ \\
$6$  & $+0.96$ & $+0.68$ \\
$8$  & $+0.97$ & $+0.76$ \\
$10$ & $+0.97$ & $+0.85$ \\
\end{tabular}
\end{ruledtabular}
\end{table}

\section{Alignment versus the held-out loss}
\label{sm:heldout}

The main text controls the apparent association between $A$ and the observable-generalization gap
$G$. The same conclusion holds when the target is the held-out loss $\Lval(\thetaK)$ itself rather
than the gap. Table~\ref{tab:heldout} reports $\rho(A,\Lval(\thetaK))$ raw, controlling for the local
improvement $\Ival$, and controlling jointly for $\Ival$ and the initial baseline $\Lval(\thetaz)$.
The raw association is moderate and negative but is removed by the joint control at every size
($|\rho|\le0.10$), so the vanishing residual is not special to the signed gap.

\begin{table}[b]
\caption{\label{tab:heldout}%
Alignment versus the held-out loss $\Lval(\thetaK)$: raw rank correlation, controlling for the local
improvement $\Ival$, and the joint control additionally removing the initial baseline
$\Lval(\thetaz)$.}
\begin{ruledtabular}
\begin{tabular}{cccc}
$n$ & $\rho(A,\Lval^{K})$ & $\rho(A,\Lval^{K}\,|\,\Ival)$ & $\rho(A,\Lval^{K}\,|\,\Ival,\Lval^{0})$ \\
\colrule
$4$  & $-0.56$ & $-0.34$ & $-0.06$ \\
$6$  & $-0.49$ & $-0.39$ & $-0.03$ \\
$8$  & $-0.45$ & $-0.38$ & $-0.10$ \\
$10$ & $-0.36$ & $-0.38$ & $-0.02$ \\
\end{tabular}
\end{ruledtabular}
\end{table}

We also report the baseline-corrected gap $\Delta G=[\Lval(\thetaK)-\Ltr(\thetaK)]-
[\Lval(\thetaz)-\Ltr(\thetaz)]=\Itr-\Ival$, which removes the initial baseline at the level of the
target. Its raw association with $A$ is $\rho(A,\Delta G)=-0.48,-0.55,-0.53,-0.59$ at $n=4,6,8,10$;
controlling for $\Ival$ alone removes it, $\rho(A,\Delta G\,|\,\Ival)=+0.03,-0.03,-0.07,-0.10$.

\section{Near-zero-mean-transfer regime}
\label{sm:transfer}

Table~\ref{tab:transfer} quantifies the transfer content of the primary task family. Support
adaptation improves the support loss substantially (mean $\Itr$ from $+0.41$ to $+0.16$), but the
mean validation improvement is statistically indistinguishable from zero at every size, with roughly
half of tasks improving. The per-task spread of $\Ival$ remains nonzero and is the quantity the
first-order identity predicts. This regime is a clean null-transfer setting by construction:
coefficients are drawn independently and the linear observable loss over disjoint Pauli support and
query subsets has no shared minimizer, so there is little mean observable transfer to detect.

\begin{table}[t]
\caption{\label{tab:transfer}%
Transfer content of the primary task family: mean and standard deviation of the validation
improvement $\Ival$, fraction of tasks with $\Ival>0$, and mean support improvement $\Itr$.}
\begin{ruledtabular}
\begin{tabular}{ccccc}
$n$ & $\langle\Ival\rangle$ & $\mathrm{std}(\Ival)$ & $\mathrm{frac}(\Ival>0)$ & $\langle\Itr\rangle$ \\
\colrule
$4$  & $-0.008$ & $0.283$ & $0.47$ & $0.410$ \\
$6$  & $+0.002$ & $0.161$ & $0.51$ & $0.313$ \\
$8$  & $-0.003$ & $0.096$ & $0.49$ & $0.242$ \\
$10$ & $+0.001$ & $0.060$ & $0.50$ & $0.158$ \\
\end{tabular}
\end{ruledtabular}
\end{table}

\section{Robustness checks}
\label{sm:robust}

The four checks below vary one ingredient of the setup at a time and confirm that the identity link
and the vanishing joint-controlled gap association both persist.

\subsection{Term-type-stratified support/query split}

To rule out that the gap is an artifact of an unbalanced split, we stratify the support/query split
so that $XX$, $YY$, $ZZ$, and $Z$ terms are balanced across the two subsets, and compare against the
uniform (random) split. Table~\ref{tab:strat} shows that both splits preserve the high identity link
and the near-zero joint-controlled gap association.

\begin{table}[b]
\caption{\label{tab:strat}%
Uniform versus term-type-stratified support/query split. The identity link $\rho(A,\Ival)$ and the
joint-controlled gap association $\rho(A,G\,|\,\Ival,\Lval(\thetaz))$ are preserved under
stratification.}
\begin{ruledtabular}
\begin{tabular}{cccc}
$n$ & split & $\rho(A,\Ival)$ & $\rho(A,G\,|\,\Ival,\Lval^{0})$ \\
\colrule
$4$ & uniform    & $0.35$ & $-0.29$ \\
$4$ & stratified & $0.50$ & $+0.10$ \\
$6$ & uniform    & $0.81$ & $-0.03$ \\
$6$ & stratified & $0.72$ & $-0.05$ \\
$8$ & uniform    & $0.76$ & $+0.04$ \\
$8$ & stratified & $0.86$ & $-0.01$ \\
\end{tabular}
\end{ruledtabular}
\end{table}

\subsection{Additional Hamiltonian families and topologies}

We repeat the analysis at $n=6$ for three Hamiltonian families---the anisotropic XYZ family of the
main text, a transverse-field Ising model (TFIM), and a fixed random-local-Pauli family---each with a
linear-chain and a ring entangler. Table~\ref{tab:families} shows that the identity link stays high
($\rho(A,\Ival)=0.49$--$0.91$) and the joint-controlled gap association stays near zero ($-0.17$ to
$+0.13$) throughout, so the mechanism is not specific to the XYZ family or the linear chain.

\begin{table}[t]
\caption{\label{tab:families}%
Identity link and joint-controlled gap association across three Hamiltonian families and two
entangler topologies ($n=6$, $L=3$).}
\begin{ruledtabular}
\begin{tabular}{llcc}
family & topology & $\rho(A,\Ival)$ & $\rho(A,G\,|\,\Ival,\Lval^{0})$ \\
\colrule
XYZ        & linear & $0.81$ & $-0.03$ \\
XYZ        & ring   & $0.89$ & $+0.08$ \\
TFIM       & linear & $0.49$ & $+0.13$ \\
TFIM       & ring   & $0.77$ & $+0.07$ \\
rand.\ Pauli & linear & $0.78$ & $-0.03$ \\
rand.\ Pauli & ring   & $0.91$ & $-0.17$ \\
\end{tabular}
\end{ruledtabular}
\end{table}

\subsection{Finite-shot variance proxy}

The observable-variance proxy $\Vtr$ is defined with a finite shot budget. To confirm that finite
sampling does not remove the effect, we compare $\Vtr$ estimated at $100$ and $1000$ shots against
the exact ($N_{\mathrm{shot}}\!\to\!\infty$) value across tasks. The rank agreement is $\rho=+0.84$
at $100$ shots and $\rho=+0.98$ at $1000$ shots, with a small positive sampling bias (mean offset
$+0.003$ and $+0.001$, respectively). Finite-shot statistics therefore add, rather than remove,
sources of the artifact.

\subsection{Random initialization}

Finally, we replace the Reptile meta-initialization with a plain random initialization
($\theta\sim\mathcal{N}(0,0.1^2)$, no meta-training) at $n=6$, $L=3$. The identity link is
$\rho(A,\Ival)=+0.74$, the raw gap association is $\rho(A,G)=-0.27$, and it collapses under the joint
control to $+0.07$---matching the Reptile arm ($0.68$, $-0.42$, $-0.01$). The identity and the
artifact are thus properties of the first-order variational geometry rather than of the
meta-initialization.

\section{Synthetic positive control: the protocol passes genuine signal}
\label{sm:positive}

Every demonstration in the main text is a collapse case, so one may ask whether the joint control is
simply too aggressive---whether it would erase a genuine alignment--gap association if one existed.
A closed-form synthetic ensemble answers this. Per task we draw an alignment $A$, a norm product $q$
with $\CV\approx0.45$ (matching the experiments), the identity-channel improvement
$\Ival=\alpha A q+\varepsilon$, and an independent baseline $\Lval(\thetaz)$; the gap is then
constructed in two ways. In the \emph{channels-only} construction,
$G=-0.8\,\Ival+0.8\,\Lval(\thetaz)+\varepsilon'$, so any $A$--$G$ association flows entirely through
the two controlled channels. In the \emph{genuine-channel} construction a direct term $-0.5\,A$ is
added, representing a real alignment$\to$gap effect beyond local improvement and baseline. Three
seeds of $150$ tasks each give Table~\ref{tab:positive}: the channels-only association collapses
under the joint control ($-0.65\to-0.06$), exactly as in the physical experiments, while the
genuine-channel association survives it ($-0.90\to-0.62$; per-seed range $[-0.66,-0.58]$). The numerical values depend on the injected channel strength: sweeping the direct coefficient from
$0.1$ to $0.8$ (channel noise s.d.\ $0.10$) gives surviving joint-controlled correlations of $-0.21$,
$-0.33$, $-0.45$, $-0.62$, and $-0.77$, so survival is robust down to injected strengths comparable
to the channel noise. The
controls of the falsification protocol therefore collapse channel artifacts without erasing genuine
signal: a diagnostic that survives them is a real candidate, not a false negative.

\begin{table}[b]
\caption{\label{tab:positive}%
Synthetic positive control. Rank correlation of alignment with the gap, raw and under the joint
control, when the gap contains only the local-improvement and baseline channels versus when a
genuine direct alignment channel is added ($3$ seeds $\times150$ tasks, pooled).}
\begin{ruledtabular}
\begin{tabular}{lccc}
construction & $\rho(A,G)$ & $\rho(A,G\,|\,\Ival)$ & $\rho(A,G\,|\,\Ival,\Lval^{0})$ \\
\colrule
channels only    & $-0.65$ & $-0.01$ & $-0.06$ \\
genuine channel  & $-0.90$ & $-0.45$ & $-0.62$ \\
\end{tabular}
\end{ruledtabular}
\end{table}

\end{document}
